%! Comparing Studies and Choosing a Method — Unit 1, Lesson 1.7 — Homework (Answer Key)
% PILOT SHAPE (2026-09-06): modeled on AP Statistics lesson 1.4 minus its AP Practice page.
% This is the lesson's SCORED individual practice and the LAST component in the packet.
% Students START IT IN CLASS in the last ten minutes of the period, alone, while the
% teacher circulates for the second formative read; they finish it at home. Packet pages are
% the default; a Desmos activity may replace them on a given day (decided at Close & Assign).
% THIRD CONTEXT: the notes teach on Riverbend, the Guided Practice runs on Harbor Point, and
% these items are on the Millbrook Public Library — recall is not the skill.
% TWO PAGES MAX (user decision 2026-09-06): two boxes — items 1-3 on page 1, items 4-6 and
% the spiralbox on page 2 — so no item breaks across a page.
% Item 4 rows (i) and (ii) are the deliberate CONTRAST PAIR: a treatment the library can hand
% out by chance against one nobody can hand out. Item 6(a) is the target misconception head-on
% (observation the method inside an experiment); 6(b) is the 48-point gap — the biggest number
% and the weakest evidence.
% THE DATA — Millbrook Public Library: 1,200 card holders; watched (1.5): 80% vs 32%, gap 48;
%   assigned (1.6): 200 volunteers 100/100, 65% vs 35%, gap 30; NEW: 300 volunteers 150/150,
%   96/150 = 64% vs 60/150 = 40%, gap 24.
% PAGE LOCKSTEP: the key is generated from this file; every work block is byte-identical.
\documentclass[12pt]{article}
\usepackage{saar-article}
\usepackage{saar-key}   % pulls in -boxes; do NOT also load -boxes
\newcommand{\TallMath}[1]{$\displaystyle #1\rule[-1.4em]{0pt}{3.2em}$}
% Word bank strip for every fill-in cluster (user request 2026-09-06). Defined per document;
% the key inherits it and prints the same strip, so heights match.
\newcommand{\wordbank}[1]{\par\vspace{2pt}\noindent\fcolorbox{goldacc}{goldbg}{\parbox{\dimexpr\linewidth-2\fboxsep-2\fboxrule\relax}{\small\textbf{Word bank:} #1}}\par\vspace{4pt}}

\begin{document}

\pageheader{Unit 1, Lesson 1.7}{Homework --- Answer Key}
% NAMESTRIP: no \namedateperiod here — mirrors the blank. NO teachernote here —
% teacher prose lives in the lesson plan.

\vspace{0.12in}

% Authored identically in homework/ and homework_key/.
\begin{remindbox}
\small
\textbf{This is your graded homework.} It is scored out of 12 and due at the start of the first
class after two study halls. You start it in the last minutes of the period --- ask about
anything that stalls you before you leave, then finish it on your own.
\end{remindbox}

\vspace{0.08in}

\boxguard[20]
\begin{scenariobox}[Millbrook Public Library --- The Third Study]{cerulean}
Millbrook has $1{,}200$ card holders and has studied its Summer Reading Program twice.
\textbf{Last summer it watched:} of the card holders who signed themselves up, $80\%$ read at
least $10$ books, against $32\%$ of those who had not --- a $48$-point gap. \textbf{This summer
it assigned:} a generator split $200$ volunteers $100$/$100$; $65\%$ against $35\%$ --- a
$30$-point gap.

\vspace{2pt}
\textbf{This fall: a third study.} Does a weekly evening book club raise how many card holders
read at least $10$ books in a semester? A generator sends $150$ of $300$ volunteers to the club
and $150$ to no club. By December, $96$ of the club group and $60$ of the no-club group had
read at least $10$ books.
\end{scenariobox}

\vspace{0.08in}

\boxguard
\begin{notesbox}{Practice}
Every answer names the design, the reason, or the method --- and the verb the study earns.
\wordbank{observational study \;\textbullet\; experiment \;\textbullet\; the card holders themselves \;\textbullet\; yes \;\textbullet\; no \;\textbullet\; impossible \;\textbullet\; unethical \;\textbullet\; impractical \;\textbullet\; survey \;\textbullet\; interview \;\textbullet\; focus group \;\textbullet\; observation \;\textbullet\; content analysis}
\begin{enumerate}[leftmargin=1.6em, itemsep=7pt, topsep=2pt]

  \item Sort the two earlier studies by design.

        \par\vspace{3pt}
        Last summer's: \ans{observational study} \hfill This summer's: \ans{experiment}
        \par\vspace{6pt}
        Who decided which group each card holder was in \emph{last} summer? \ans{the card holders}

  \item Compute the book-club rates, then find the gap.

        \begin{work}
          \dfrac{96}{150} &= 0.64 = 64\% \\
          \dfrac{60}{150} &= 0.40 = 40\%
        \end{work}

        The club group is higher by \ans{24} percentage points. Write the strongest
        sentence this study earns.
        \par\ansline{The weekly book club caused more card holders to read at least 10 books ---}
\ansline{a generator built the two groups, so the club is the only difference.}

\end{enumerate}
\end{notesbox}

\boxguard[30]
\begin{notesbox}{Practice, continued}
\begin{enumerate}[leftmargin=1.6em, itemsep=5pt, topsep=2pt, start=3]

  \item Fill in the plan the director followed for the book-club study.

        \nopagebreak
        \renewcommand{\arraystretch}{1.15}
        \begin{tabularx}{\linewidth}{@{} >{\raggedright\arraybackslash}p{3.5cm} X @{}}
        \toprule
        \textbf{Stage} & \textbf{The book-club study} \\
        \midrule
        Formulate & Does the club raise how many card holders read $10$ books? \\
        Collect & \ans{Generator splits 300 volunteers 150/150.} \\
        Represent & \ans{A table: each group's count and percent.} \\
        Analyze \& report & \ans{Compare 64\% vs.\ 40\%; say it in context.} \\
        \bottomrule
        \end{tabularx}


  \item The board sends two more questions. For each, decide whether the director could
        assign the groups --- and if not, say which reason.

        \nopagebreak
        \renewcommand{\arraystretch}{1.15}
        \begin{tabularx}{\linewidth}{@{} X p{1.4cm} p{3.6cm} @{}}
        \toprule
        \textbf{The question} & \textbf{Assign?} & \textbf{If not, why?} \\
        \midrule
        Does a $\$5$ late fee reduce overdue books? & \ans{yes} & \ans{--- experiment} \\
        Do children read to as toddlers read more at $12$? & \ans{no} & \ans{impossible} \\
        \bottomrule
        \end{tabularx}

  \item Name the collection method that fits each job best.

        \nopagebreak
        \renewcommand{\arraystretch}{1.15}
        \begin{tabularx}{\linewidth}{@{} X p{4.2cm} @{}}
        \toprule
        \textbf{What the library needs} & \textbf{Best method} \\
        \midrule
        Which titles were borrowed most over the past five years & \ans{content analysis} \\
        How many of the $1{,}200$ card holders want Sunday hours & \ans{survey} \\
        Why one long-time member stopped coming in, in her own words & \ans{interview} \\
        How many visitors actually use the new self-checkout machine & \ans{observation} \\
        Six teens talking together about what the teen room needs & \ans{focus group} \\
        \bottomrule
        \end{tabularx}

  \item Two claims about the three studies.

        \textbf{(a)} She collected the book-club data by sitting in on the meetings and
        recording who actually came, so a classmate calls it an observational study. Name the
        method and the design, and say whether the classmate is right.
        \par\ansline{Method: observation. Design: experiment --- a generator built the 150/150 groups.}
\ansline{The classmate is wrong: how the data were recorded never decides the design.}

        \textbf{(b)} Millbrook now has a $48$-point gap, a $30$-point gap, and a $24$-point
        gap. Which is the \emph{weakest} evidence that the programs work --- and why is it the
        one with the biggest number?
        \par\ansline{The 48-point gap --- the card holders who signed themselves up were already the}
\ansline{heavy readers, so the gap measures who signed up, not what the program did.}

\end{enumerate}
\end{notesbox}

\boxguard[3]   % the spiralbox needs about three baselines; a larger guard pushes it to a third page
\begin{spiralbox}
\textbf{Next lesson: Lesson 1.8 --- You Run the Study.} From a question your class chooses,
you will decide whether it can be an experiment at all, plan the cycle, pick a sampling
technique and a collection method, and report a conclusion your design earns.
\end{spiralbox}

\end{document}
